\scanpage{47}
\mspage{6}
\uline{Corollaire} 1 Soient $\unsure{G_1|G}$ $\deleted{\unsure{G_2|G}}\unsure{H|H_2}$ deux \deleted{\unsure{sous}} schémas en groupes comm. de type fini sur $S$ loc. noeth. connexe, $\unsure{G_2|G}$ \uline{rigide}. Soient $u,v:G\longrightarrow H$ \unsure{des|deux} hom de $S$-groupes, et $s\in S$ tel que $u_s,v_s:G_s\rightrightarrows H_s$ soient égaux. Alors $u=v$

On applique le th. de rigidité \unsure{aux|au} graphes

\uline{Corollaire 2} Soient $S$ un schéma loc. noeth, $S'\longrightarrow S$ un morphisme \uline{fid plat} \deleted{\unsure{et quasi-cpt}}\inserted{\uline{de type fini}} avec $S'$ loc noeth, $G$ un préschéma en groupes comm de type fini sur $S$, $H'$ un sous-schéma en groupes fermé de $G'=G\times_S S'$. On suppose

(i) les fibres de $S'\unsure{\longrightarrow|\text{---}}S$ sont \unsure{abslt.|abst.|géom.} connexes

(ii) $H'$ est rigide

Alors il existe un sous-schéma en groupes \inserted{fermé} $H$ de $G$ (nécessairement unique, et rigide) tel que $H'=H\times_S S'$.

On applique la théorie de la descente, pour laquelle il faut seulement
% archive sheet number: 46
